\subsection{Lenguajes y Herramientas [Pendientes: 1. Elegir tecnologías de aplicación y pruebas y proveedor de API. 2. Fijar versiones. 3. Configurar repositorio y reglas de Git. 4. Confirmar acceso documental definitivo y recursos HPC. 5. Buscar y recopilar referencias bibliográficas oficiales, manuales de usuario y documentación técnica de todas las herramientas y lenguajes utilizados, incluidos FDS, Smokeview, UML, Git, GitHub, Aldine, LaTeX, Docker Compose, Nginx, Tailscale, Linux y Slurm.]}\label{sec:spmp-8-2}
\begin{enumerate}
\item \textbf{Simulación.} FDS es el software utilizado para preparar los escenarios de incendio mediante simulaciones, según la convención FDS/WFDS de la sección \ref{sec:spmp-6-1}. La documentación oficial \textit{Installing and Running FDS on a Linux Cluster} describe su instalación y ejecución en clústeres Linux, incluye ejemplos con el gestor de trabajos Slurm y permite acceder al repositorio oficial del software~\cite{fdslinuxcluster}. Esta base tecnológica es acorde con el entorno Linux, Slurm y la arquitectura x86 de la infraestructura HPC-Zine prevista para el proyecto; la correspondencia general entre estas tecnologías se distingue de la comprobación de su instalación, configuración y ejecución efectiva en ese recurso.

\item \textbf{Visualización.} Smokeview es el visualizador previsto para presentar los resultados de FDS mediante las animaciones georreferenciadas de los escenarios.

\item \textbf{Modelado.} UML 2.5.1 proporciona el lenguaje para representar y comunicar el diseño del software~\cite{uml}.

\item \textbf{Control de versiones y conservación.} Git conserva las versiones del código de la aplicación y los repositorios documentales por proyecto.

\item \textbf{Documentación.} Aldine centraliza la elaboración colaborativa de los documentos y la bibliografía del proyecto; LaTeX es el lenguaje empleado para estructurar esos documentos.

\item \textbf{Infraestructura y administración.} Docker Compose organiza los contenedores de aplicación y compilador del servidor documental; Nginx actúa como gateway y Tailscale permite el acceso al servidor documental entre los miembros del equipo y el director del proyecto. Este computador está ubicado en casa de Daniel y dispone de almacenamiento persistente. Para la ejecución científica, HPC-Zine es la infraestructura prevista, Linux es el entorno de ejecución y Slurm gestiona los trabajos de simulación. Los ámbitos de infraestructura se describen en la sección \ref{sec:spmp-9-2-2}~\cite{fdslinuxcluster}.
\end{enumerate}

\subsubsection{Análisis de Alternativas y Justificación [Pendientes: 1. documentar comparación de simuladores y herramientas aún por elegir]}\label{sec:spmp-8-2-1}
WFDS se mantiene como la elección que el equipo considera más adecuada para el tipo de representación física y de escenarios que pretende preparar. La decisión vigente es utilizarlo de forma exclusiva, mientras que FARSITE se conserva como antecedente del proceso de formulación. Esta justificación debe vincularse con las necesidades y restricciones de las secciones \ref{sec:spmp-6-2} y \ref{sec:spmp-6-3}, y completarse con la evidencia técnica y las comparaciones ya disponibles, porque la preferencia del equipo no constituye una demostración de que sea el mejor simulador para todos los contextos. La ejecución del piloto será una comprobación de viabilidad de la elección y no una comparación retrospectiva ficticia con herramientas que no se hayan ensayado~\cite{diseno2026}.

UML se elige porque aporta convenciones compartidas para comunicar el diseño y porque sus conceptos son conocidos por el equipo como estudiantes de Ingeniería de Sistemas de la Pontificia Universidad Javeriana. Aldine se adopta como reemplazo de Overleaf para disponer de un espacio centralizado de documentos sin depender del pago de uso de esa plataforma, aprovechando el servidor disponible. Esta motivación no presupone que el computador, su operación o sus respaldos carezcan de costos. Git, por su parte, permitirá conservar la evolución del código de la aplicación y relacionar cambios con versiones comprobadas. Las decisiones anteriores ofrecen razones de selección, aunque la comparación de otras tecnologías de desarrollo, pruebas y almacenamiento todavía requiere fuentes y criterios documentados~\cite{uml,servidor2026,adaptacion2026}.

